
\begin{frame}
{Oracle}

Principales structures physiques dans ORACLE\,:
\begin{redbullet}
  \item  {\bf bloc}: unité physique d'E/S.\\
     La taille d'un bloc ORACLE est un
        multiple de la taille des blocs du système sous-jacent.

  \item {\bf Extension}:  ensemble de blocs contigus contenant
        un m\^eme type d'information.

  \item  {\bf Segment}: ensemble d'extensions stockant un 
        objet logique (une table, un index ...). 
\end{redbullet}

Le paramétrage du stockage des données (tables et index)
est spécifié dans un {\it tablespace}

\end{frame}


\begin{frame}{Tables, segments, extensions et blocs}

\figSlideWithSize{oraphys}{8cm}

\end{frame}

\begin{frame}{Les blocs Oracle}

La structure d'un bloc repose sur
un adressage physique/logique, chaque enregistrement
ayant une adresse interne.

\vfill
 
\figSlideWithSize{bloc-oracle}{8cm}

\vfill

Un chaînage est créé quand il faut déplacer un enregistrement.

\end{frame}

\begin{frame}{Gestion de l'espace}

\begin{redbullet}
 \item \sqltxt{PCTFREE} donne l'espace libre à préserver 
  au moment de la création d'une table ou d'un index.

  \item  \sqltxt{PCTUSED} indique le seuil en dessous duquel  le bloc
              est disponible pour des insertions.
\end{redbullet}

Oracle maintient un répertoire (?) des blocs
disponibles pour insertions.

\vfill
 Exemple\,:

\begin{redbullet}
   \item  \sqltxt{PCTFREE} = 30\% et  \sqltxt{PCTUSED}=70\%
   \item  \sqltxt{PCTFREE} = 10\% et  \sqltxt{PCTUSED}=80\%
\end{redbullet}

Le second choix est plus efficace, mais plus risqué et plus coûteux.

\end{frame}


\begin{frame}{Extensions et segments}

\alert{L'extension} est une suite de blocs contigus.

\vfill

Le \alert{segment} est un ensemble d'extensions
 contenant un objet logique. 

Il existe quatre types de segments\,:
\begin{redbullet}
  \item Le segment de données.
  \item Le segment d'index.
  \item Le $rollback\ segment$ utilisé pour les transactions.
  \item Le segment temporaire (utilisé pour les tris).
\end{redbullet}

\vfill

Moins il y a d'extensions dans un segment, plus il
est efficace.

\end{frame}


\begin{frame}{Stockage et adressage des enregistrements}

En règle générale un 
enregistrement est stocké dans un seul bloc.
\vfill


L'adresse physique d'un enregistrement est le $ROWID$. Il contient:
\vfill

\begin{redbullet}
   \item Le numéro du bloc dans le {\bf fichier}.
   \item Le numéro du n-uplet dans lebloc.
   \item le numéro du segment.
   \item Le numéro du fichier.
\end{redbullet}
Exemple\,: 00000DD5.000.2.001 est l'adresse du premier n-uplet
du bloc DD5 dans le second segment du premier fichier.

\end{frame}


\begin{frame}{Les {\it tablespaces}}

Une base est divisée par l'administrateur
   en $tablespace$. Chaque $tablespace$ consiste en un (au moins)
        ou plusieurs fichiers.

\vfill

La notion de $tablespace$ permet\,:

\vfill

\begin{redbullet}
  \item De contr\^oler l'emplacement physique des données.
        (par ex.\,: le dictionnaire sur un disque,
        les données utilisateur sur un autre).

  \item de régler l'allocation de l'espace (extensions).

  \item de faciliter la gestion (sauvegarde, protection, etc).
\end{redbullet}

\end{frame}

\begin{frame}{Exemple de  {\it tablespaces}}

\figSlideWithSize{oratbspace}{8cm}

\end{frame}

\begin{frame}[fragile]{Création de  {\it tablespaces}}

\begin{minted}[frame=leftline,
               baselinestretch=.8,
               fontsize=\footnotesize,
               framesep=2mm]{sql}
CREATE TABLESPACE TB1
    DATAFILE 'fichierTB1.dat' SIZE 50M
      DEFAULT STORAGE (
        INITIAL 100K, NEXT 40K, MAXEXTENTS 20,
         PCTINCREASE 20);
\end{minted}

\vfill
\begin{minted}[frame=leftline,
               baselinestretch=.8,
               fontsize=\footnotesize,
               framesep=2mm]{sql}
CREATE TABLESPACE TB2 
    DATAFILE 'fichierTB2.dat' SIZE 2M
    AUTOEXTEND ON NEXT 5M MAXSIZE 500M
    DEFAULT STORAGE (INITIAL 128K
                       NEXT 128K MAXEXTENTS UNLIMITED);
\end{minted}


\end{frame}

\begin{frame}{Maintenance d'un \tbspace}

Quelques actions 
disponibles sur un \tbspace\,:

\begin{redbullet}
  \item On peut mettre un \tbspace\ hors-service.

\texttt{ALTER TABLESPACE TB1 OFFLINE;}

  \item On peut mettre un \tbspace\ en lecture seule.

\texttt{ALTER TABLESPACE TB1 READ ONLY;}

  \item   On peut ajouter un nouveau fichier.

\texttt{ALTER TABLESPACE ADD DATAFILE \\
'fichierTB1-2.dat' SIZE 300 M;}

\end{redbullet}

\end{frame}

\begin{frame}[fragile]{Affectation de tables à un \tbspace}

On peut placer une table dans un \tbspace. Elle
prend alors les paramètres de stockage de ce dernier.

\vfill

On peut aussi remplacer certaines valeurs.

\vfill

\begin{minted}{sql}
create table Film (...)
pctfree 10
pctused 40
tablespace TB1
storage ( initial 50K
          next 50K
         maxextents 10
pctincrease 25 );
\end{minted}

\end{frame}